\documentclass[a4paper]{article}
\usepackage[margin=1in]{geometry}
\usepackage{xcolor, colortbl}
\usepackage{hyperref}
\usepackage{graphicx}
\usepackage{amsmath}
\usepackage{amsfonts}
\usepackage{amssymb}
\usepackage{multicol}
\usepackage{listings}
\usepackage{subcaption}
\usepackage{microtype}

\definecolor{codebg}{HTML}{F7F7F2}
\definecolor{codekeyword}{HTML}{1D4E89}
\definecolor{codecomment}{HTML}{5E6B5E}
\definecolor{codestring}{HTML}{8A4F1D}
\definecolor{codetype}{HTML}{7B2CBF}

\lstdefinelanguage{AFPLang}{
    sensitive=true,
    morekeywords={let,mut,fn,return,if,then,else,while,spawn,yield,push,insert,remove,module,struct,end,type,fun,match,with,open,rec,in,and},
    morekeywords=[2]{int,bool,unit,True,False,Some,None,Ok,Error,StateT,Either,Result,Deferred},
    morecomment=[l]{--},
    morecomment=[s]{(*}{*)},
    morestring=[b]"
}

\lstset{
    language=AFPLang,
    frame=single,
    framesep=5pt,
    backgroundcolor=\color{codebg},
    basicstyle=\small\ttfamily,
    columns=fullflexible,
    keepspaces=true,
    breaklines=true,
    tabsize=2,
    inputencoding=utf8,
    extendedchars=true,
    showstringspaces=false,
    keywordstyle=\color{codekeyword}\bfseries,
    keywordstyle=[2]\color{codetype},
    commentstyle=\color{codecomment}\itshape,
    stringstyle=\color{codestring}
}

\hypersetup{colorlinks=true, linkcolor=blue!60!black, urlcolor=blue!60!black}
\setcounter{secnumdepth}{3}
\setlength{\parindent}{0pt}
\setlength{\parskip}{\baselineskip}
\renewcommand{\thesubsection}{\arabic{subsection}}
\renewcommand{\thesubsubsection}{\thesubsection.\arabic{subsubsection}}

\title{\textbf{Ownership Types}\\}
\author{Alexandru Bolfa - 5495741}

\begin{document}

\maketitle

\subsection*{Introduction}

This report showcases the \textbf{implemented features} of the ownership
language and explains how the project meets the \textbf{one-person AFP
requirements}. It covers the syntax, the implemented feature phases, example
programs, and the required functional-programming techniques used in the
implementation: \textbf{persistent data structures}, \textbf{Accessor-based
lenses/traversals}, and \textbf{monad transformers/effects}.

The project README\footnote{The project README is also included as
Appendix~\ref{appendix:readme}.} is the practical guide for building and running
the project. This report is the \emph{design and architecture overview:} what
was built, how the type checker and interpreter are organized, where the
required techniques appear in the code, and why the \textbf{persistent borrow
graph} counts as advanced usage.

Section 1 outlines the \emph{syntax}, \emph{implemented features}, and
\emph{pass pipeline}. Section 2 covers the \textbf{required techniques}. Section
3 describes the persistent borrow graph used for \textbf{advanced usage}.
Section 4 summarizes the type and ownership checks behind the implemented
features.

\subsection{Design and Scope}

\subsubsection*{Syntax}

The language is a \emph{small Rust-like subset}. It has ownership, borrowing,
mutation, lists, functions, and cooperative concurrency:

\begin{lstlisting}
R ::= int | bool | unit | [R]
T ::= R | &R | &mut R

S ::= let x = E | let mut x = E | P = E
    | fn f(params) -> T { S* } | { S* }
    | if E then { S* } | if E then { S* } else { S* }
    | while E { S* } | return E | E

E ::= n | True | False | x | E op E | if E then E else E
    | f(E*) | [E*] | E[E] | &P | &mut P | *E
    | spawn(f(E*)) | yield()

P ::= x | P[E] | *P
\end{lstlisting}

\texttt{unit} is the type for computations that do \emph{not} produce useful data.
It has one runtime value, written as \texttt{()}, and is used by effect-only
operations such as \texttt{yield()}, \texttt{spawn(f(args))}, and list mutation.

\subsubsection*{Implemented features:}

\begin{center}
\begin{tabular}{r | p{.68\linewidth}}
Phase 0 & Mutable variables, blocks and control flow, and more complex
functions. \\
Phase 1 & Affine types. \\
2A & Lists. \\
2B & Copyable primitives. \\
3A & Immutable borrows. \\
3B & Mutable borrows. \\
4B & Multi-threading through \texttt{spawn} and \texttt{yield}. \\
\end{tabular}
\end{center}

These features cover the \textbf{one-person phase targets}: phase 0 and phase
1, two phase-2 features (\texttt{2A}, \texttt{2B}), two phase-3 features
(\texttt{3A}, \texttt{3B}), and one phase-4 feature (\texttt{4B}). Together
with the \textbf{advanced persistent borrow graph} in Section 3, this covers the
feature and technique requirements listed for a one-person submission.

\subsubsection*{Pipeline}

The implementation is organized as a \textbf{small pass pipeline}. Expressions
and statements are kept separate in Core, and assignments/borrows use a shared
\emph{place} representation such as \texttt{x}, \texttt{xs[0]}, or \texttt{*r}.

\begin{center}
\small
\setlength{\fboxsep}{6pt}
\begin{tabular}{ccccccc}
\fbox{\strut Lowering} & $\rightarrow$ &
\fbox{\strut Typecheck} & $\rightarrow$ &
\fbox{\strut Borrowcheck} & $\rightarrow$ &
\fbox{\strut Interpret}
\end{tabular}
\end{center}

\textbf{Lowering} gives names unique identities, \textbf{typechecking} assigns
ordinary types, \textbf{borrow checking} enforces moves and active borrows, and
\textbf{interpretation} runs only after both static checks pass. Runtime
references are symbolic paths into the environment, so writing through
\texttt{\&mut xs[0]} rebuilds the path back to \texttt{xs[0]}.

\subsubsection*{Tools and Libraries}

The corresponding \texttt{.mli} files document the higher-level functions and
state types exposed by each pass. The tests include \textbf{positive examples}
that should run and \textbf{negative examples} that should be rejected by name
resolution, typechecking, borrow checking, or runtime checks.

The implementation deliverables are the \textbf{typechecker},
\textbf{interpreter}, \textbf{positive and negative test suite},
\textbf{module documentation}, and this design overview.

The main libraries are \texttt{Base} for data structures
\cite{janestreet_base}, \texttt{Async} for \texttt{spawn}/\texttt{yield}
\cite{janestreet_async}, \texttt{Accessor} for generated record accessors
\cite{janestreet_accessor}, \texttt{monadic} for monad transformers
\cite{denommus_monadic}, and \texttt{dune}/\texttt{ppx\_jane} for building and
PPX support \cite{dune,janestreet_ppx_jane}.

\subsection{Implementation Techniques}

For a one-person submission, the required techniques are \textbf{persistent
functional data structures}, \textbf{Accessor-based lenses}, and \textbf{monad
transformers}.

\subsubsection{Persistent Functional Data Structures}

Environments are \emph{immutable} \texttt{Map.M(String)} values wrapped by
\texttt{Common.Value.Env}. Updating a variable or function binding returns a new
environment. Lists are also \textbf{persistent}: \path{src/common/plist.ml}
rebuilds only the path that changes, which matters for writes through places such as
\texttt{\&mut xs[0]}.

In Haskell terms, the environments are used like persistent \texttt{Map} values:
an update does not mutate the old environment, but returns a new version that
shares unchanged structure with the previous one. The OCaml code follows the
same model with \texttt{Base.Map}. For example, binding a variable uses
\texttt{Env.bind}, which returns a new map-backed environment instead of editing
the old one in place.

The custom persistent list and borrow graph use the same idea for project-specific
data. A list update rebuilds only the affected path, and the borrow graph
operations such as \texttt{add\_edge}, \texttt{remove\_owner}, and
\texttt{merge} all return new graph values. This is the same style as writing
small immutable tree or graph transformations in Haskell.

\subsubsection{Lenses / Traversals}

Pass states are records with \textbf{generated accessors}. The local monadic
wrappers expose \texttt{access} and \texttt{map\_access}, so a rule updates
only the field it needs:

\begin{lstlisting}
M.map_access owns ~f:(Env.bind x Moved)
\end{lstlisting}

The OCaml mechanism is \texttt{ppx\_accessor} plus Jane Street's
\texttt{Accessor} library. Each pass state is a record marked with
\texttt{[@@deriving accessors]}, which generates values such as \texttt{owns},
\texttt{graph}, \texttt{vars}, and \texttt{funs}. These generated accessors play
the same role as Haskell lenses for record fields.

The helper \texttt{M.access} corresponds to a Haskell-style \texttt{view}, and
\texttt{M.map\_access} corresponds to applying an \texttt{over}-style update
inside the current state:

\begin{lstlisting}
-- Haskell shape
view owns state
modify (over owns (Env.bind x Moved))
\end{lstlisting}

The OCaml code spells this through modules and functions instead of Haskell type
classes:

\begin{lstlisting}
let access accessor =
  AppM.create (fun s -> Except.pure (Accessor.get accessor s, s))

let map_access accessor ~f =
  AppM.create (fun s -> Except.pure ((), Accessor.map accessor s ~f))
\end{lstlisting}

This is used throughout the passes, not just once. For example, the borrow
checker updates \texttt{owns} and \texttt{graph} independently, while the
typechecker updates \texttt{vars}, \texttt{muts}, \texttt{funs}, and
\texttt{ret} through the same accessor pattern.

\subsubsection{Monad Transformers / Effects}

The static passes use the same structure as Haskell's \textbf{\texttt{StateT
AppState (Either String)}}, built with \texttt{monadic}:

\[
  \alpha\ AppM.t = AppState.t \to ((\alpha \times AppState.t),\ string)\ result.
\]

The main difference is syntactic. Haskell can write a transformer stack directly
as a type alias, while the OCaml library builds the same stack through modules
and functors. The Haskell shape is:

\begin{lstlisting}
type AppM a = StateT AppState (Either String) a
\end{lstlisting}

The OCaml version first wraps the state type in a module, then creates a
\texttt{Result} monad whose error type is \texttt{string}, and finally places
state on top of that result layer:

\begin{lstlisting}
module AppState = struct
  type t = env
end

module Except = Monadic.Result.Make (struct
  type t = string
end)

module AppM = Monadic.State.MakeT (Except) (AppState)
\end{lstlisting}

This gives the pass \texttt{AppM.get} and \texttt{AppM.put}, corresponding to
Haskell's \texttt{get} and \texttt{put} from the state layer. It also gives
\texttt{AppM.elevate}, which is the local equivalent of Haskell's \texttt{lift}.
For example, throwing an error from the result layer is:

\begin{lstlisting}
let fail msg = AppM.elevate (Except.error msg)
\end{lstlisting}

which corresponds to the Haskell expression:

\begin{lstlisting}
fail msg = lift (Left msg)
\end{lstlisting}

Each static pass then wraps this in a small local \texttt{M} module so the rest
of the pass can use \texttt{let\%bind}, \texttt{get}, \texttt{put},
\texttt{fail}, \texttt{access}, and \texttt{map\_access} without repeating the
transformer plumbing.

The interpreter puts \textbf{\texttt{Async.Deferred.t}} under the same stack so
\texttt{spawn} and \texttt{yield()} can suspend cooperatively. Conceptually,
that interpreter stack is the OCaml equivalent of putting an asynchronous base
monad under \texttt{ExceptT String} and then putting \texttt{StateT Env} on top:

\begin{lstlisting}
type EvalM a = StateT Env (ExceptT String Deferred) a
\end{lstlisting}

In the OCaml code, \texttt{Result.MakeT} builds the result transformer over
\texttt{Deferred}, and \texttt{State.MakeT} adds the state layer above it. The
helper \texttt{lift\_deferred} lifts scheduler actions such as \texttt{yield()}
from \texttt{Deferred.t} into the interpreter monad.

\subsection{Advanced Usage}

For \textbf{advanced usage}, the borrow checker uses a \textbf{persistent borrow
graph} implemented in \path{src/common/borrow_graph.ml}. The graph is carried
in the state of \path{src/passes/borrowcheck.ml}, so each borrow-checking step
receives one version of the graph and returns the next version.

The graph stores borrow relationships as \emph{immutable edges}:

\begin{lstlisting}
type owner =
  | Stored of AbsLang.ident
  | Temporary

type edge =
  { owner : owner
  ; target : target
  }
\end{lstlisting}

A target is the borrowed root variable together with the borrow kind
(\texttt{\&} or \texttt{\&mut}). A \emph{stored edge} means ``reference variable
\texttt{r} keeps root \texttt{x} borrowed'', and a \emph{temporary edge} means
``this expression borrows \texttt{x} until the current statement ends''.
Adding, removing, merging, or clearing edges returns a \textbf{new graph}.

The core operations are small, but they are used throughout the borrow checker:

\begin{lstlisting}
let add_edge edge graph =
  if List.exists graph ~f:(same_edge edge)
  then graph
  else edge :: graph

let add_owned ~owner target graph =
  add_edge { owner = Stored owner; target } graph

let merge a b =
  List.fold b ~init:a ~f:(fun graph edge -> add_edge edge graph)
\end{lstlisting}

The graph is part of the \textbf{borrow-checker state}:

\begin{lstlisting}
type env =
  { owns : ownership Env.t
  ; graph : Borrow_graph.t
  }

let map_graph f = M.map_access graph ~f
\end{lstlisting}

\textbf{Borrow-graph updates} go through this state field. The checker calls
\texttt{Borrow\_graph.active} before reads,
moves, assignments, and borrows to find active edges for a root. It calls
\texttt{add\_owned\_all} when a reference variable is bound, \texttt{remove\_owner}
when a reference goes out of scope or is overwritten, and \texttt{add\_temporary}
when an expression creates a short-lived borrow. Control flow also uses the same
structure: blocks call \texttt{keep\_outer}, \texttt{if} expressions and
statements call \texttt{merge} or \texttt{merge\_stored}, and each statement
finishes with \texttt{clear\_temporaries}. The important property is that these
operations \emph{return new graphs instead of mutating the previous one}.

\subsection{Type and Ownership Checks}

This section summarizes the checks that support the implemented features.

\subsubsection{Typechecking}

\textbf{Lowering} resolves source names to unique identifiers before
typechecking. The typechecker then records variable types, binding mutability,
and function signatures. It checks the ordinary shape of expressions: literals,
lists, indexing, conditionals, function calls, assignments, borrows, and
dereferences.

Assignments are checked against the type of the target \emph{place}. The place
must be writable, so assigning to an immutable binding is rejected before the
borrow checker runs. Borrowing is also typed here: \texttt{\&p} produces a
\emph{shared reference}, \texttt{\&mut p} requires a writable place and produces
a \emph{mutable reference}, and dereferencing requires either kind of reference.

\subsubsection{Ownership and Borrows}

Borrow checking runs \emph{after} typechecking. It tracks whether each owned value is
still available or has been moved, and it stores active borrow edges in the
persistent graph described in Section 3. The copyable values are
\texttt{int}, \texttt{bool}, \texttt{unit}, and shared references. Lists and
mutable references are \textbf{affine}: moving one transfers ownership.

\textbf{Reads, moves, assignments, and borrows} all consult the active borrow
graph. A move, assignment, or mutable borrow requires the root variable to have
no active borrows. A shared borrow is allowed unless there is already an active
mutable borrow. Temporary borrows are cleared at the end of the statement, while
stored borrows stay attached to the reference variable that owns them.

\subsubsection{Spawning and Runtime Checks}

\texttt{spawn(f(args))} is accepted only for function calls whose argument types
contain no references. \texttt{yield()} has type \texttt{unit} and is handled by
the interpreter's Async-based scheduler. List bounds are checked by the
interpreter during execution.

\begin{thebibliography}{1}

\bibitem{denommus_monadic}
{Denommus}.
\newblock Monadic: A monad library for OCaml.
\newblock \url{https://github.com/Denommus/monadic}, 2024.

\bibitem{janestreet_accessor}
{Jane Street}.
\newblock Accessor: Nested functional data structure accessors.
\newblock \url{https://github.com/janestreet/accessor}, 2026.

\bibitem{janestreet_async}
{Jane Street}.
\newblock Async: Jane Street's asynchronous execution library.
\newblock \url{https://github.com/janestreet/async}, 2026.

\bibitem{janestreet_base}
{Jane Street}.
\newblock Base: Standard library for OCaml.
\newblock \url{https://github.com/janestreet/base}, 2026.

\bibitem{janestreet_ppx_jane}
{Jane Street}.
\newblock ppx\_jane: Jane Street PPX rewriters.
\newblock \url{https://github.com/janestreet/ppx_jane}, 2026.

\bibitem{dune}
{OCaml Dune Project}.
\newblock Dune: Composable build system for OCaml.
\newblock \url{https://github.com/ocaml/dune}, 2026.

\end{thebibliography}

\clearpage
\appendix
\setcounter{subsection}{0}
\setcounter{subsubsection}{0}
\renewcommand{\thesubsection}{\Alph{subsection}}
\renewcommand{\thesubsubsection}{\thesubsection.\arabic{subsubsection}}
\begin{center}
  {\Huge\bfseries Appendix}
\end{center}
\vspace{1em}

\subsection{Example Programs}
\label{appendix:examples}

These are the small programs used to show the implemented feature groups. They
can be run with \texttt{make run FILE=examples/<name>.lang}.

\begingroup
\lstset{basicstyle=\scriptsize\ttfamily, framesep=3pt}
\noindent\begin{minipage}[t]{0.48\linewidth}

\subsubsection{Phase 0: functions and control flow}
\label{appendix:example-phase0}

\texttt{examples/phase0.lang} returns \texttt{120}.
\begin{lstlisting}
fn fact (n : int) -> int {
  if (n == 0) then {
    return 1;
  };
  return n * fact(n - 1)
};
return fact(5);
\end{lstlisting}

\subsubsection{Phase 1: affine ownership}
\label{appendix:example-phase1}

\texttt{examples/phase1.lang} returns \texttt{0}.
\begin{lstlisting}
fn first (xs : [int]) -> int {
  return xs[0]
};

let xs = [10, 20];
let ys = xs;

first(ys);

return 0;
\end{lstlisting}

\subsubsection{2A/2B: lists and copyable primitives}
\label{appendix:example-phase2}

\texttt{examples/phase2.lang} returns \texttt{[4, 2, 3]}.
\begin{lstlisting}
let x = 1;
let y = x;
let mut xs = [x, y];

push(xs, 3);
insert(xs, 1, 4);
remove(xs, 0);

xs[1] = xs[1] + y;

return xs
\end{lstlisting}

\end{minipage}\hfill
\begin{minipage}[t]{0.48\linewidth}

\subsubsection{3A: immutable borrows}
\label{appendix:example-phase3a}

\texttt{examples/phase3a.lang} returns \texttt{42}.
\begin{lstlisting}
fn read (x : &int) -> int {
  return *x
};

let value = 21;
let a = &value;
let b = &value;

return read(a) + read(b);
\end{lstlisting}

\subsubsection{3B: mutable borrows}
\label{appendix:example-phase3b}

\texttt{examples/phase3b.lang} returns \texttt{42}.
\begin{lstlisting}
fn set (x : &mut int) -> int {
  *x = 42;
  return 0
};

let mut value = 1;

set(&mut value);

return value;
\end{lstlisting}

\subsubsection{4B: thread-safe spawn and yield}
\label{appendix:example-phase4b}

\texttt{examples/phase4b.lang} spawns two finite cooperative tasks.
\begin{lstlisting}
fn worker (mut n : int) -> unit {
  while n > 0 {
    yield();
    n = n - 1
  };
  return yield()
};

spawn(worker(2));
spawn(worker(3));
return yield()
\end{lstlisting}

\end{minipage}
\endgroup

\clearpage

\subsection{Project README}
\label{appendix:readme}

The project README is included here for completeness. It is the practical guide
for installing, testing, and running the implementation.

\begingroup
\lstset{basicstyle=\small\ttfamily, framesep=5pt}
\begin{lstlisting}
# ownership-interp

An OCaml interpreter for a small Rust-like language with ownership types.

This README is for installing and running the project. The complementary report
is at `report/main.pdf`. It explains the implementation choices and how the
project maps to the grading criteria.

## Install

The project needs OCaml 5.2 or newer. For a fresh local switch:

```sh
opam switch create . ocaml-base-compiler.5.2.0 -y
eval $(opam env)
opam pin add monadic https://github.com/Denommus/monadic.git#42c0a97c0667ee3daec8a31b26314b696243ee8d -y
opam install . --deps-only --with-test -y
```

## Run

```sh
make build
make test
make run                           # REPL
make run FILE=examples/phase0.lang # run a file
```

Example programs live in `examples/`.
\end{lstlisting}
\endgroup

\end{document}
